\section{Introduction}

Extant consensus protocols fail to provide the necessary performance to substantiate their utility, with the limiting factors of these protocols exhibited by high operation costs or limited performance outcomes. While new protocols continue to be developed which are progressively faster than their predecessors, the scalability of said protocols remains non-existent, with performance outcomes remaining constant despite network size. In developing a protocol that displays constant performance, every additional node is a burden on the precursory members. For such protocols, there will always exist an incentive to reduce network participation, which in turn has negative consequences regarding a platform's decentralization. Given these algorithms are primarily designed for substantiating guaranteed network performance in a trustless environment, increased decentralization is considered a necessity.
A majority of these protocols present additional risks, many of which have not been well understood prior to their real-world implementation. For the sake of brevity, we will be limiting the examples of such failings to the Bitcoin network and protocol, with the understanding that many other systems exhibit similar failings as a consequence of their consanguineous nature \cite{Nakamoto2008}. As to be discussed in greater detail later in this paper, the failings include the implementation of competitive proofs, the consolidation of resources (compute systems or currency) outside the domain of the decentralized network and the non-deterministic nature of validation proofs.

Working to mediate the aforementioned failings and limitations of existing protocols, this paper presents a novel consensus protocol, \texttt{Blossom}, with a variable tolerance for Byzantine (i.e., arbitrary) faults (\(f \approx \frac{1}{3}\)), deterministic transaction finality and scalable performance with difficulty increases at (non-constant) logarithmic time. The result is a protocol that incentivizes network scaling, introduces supermajority consensus requirements, and mitigates the risks associated with resource accumulation. 
The most salient risk to our protocol, a Denial-of-Service attack (DoS) was mitigated by removing synchronicity requirements and implementing redundancy measures to guarantee the continuity of consensus despite a node faltering mid-committee. Other risks which we have mitigated in our protocol include Sybil attacks, Routing attacks, and Distributed-Denial-of-Service (DDoS) attacks. In implementing our protocol we recommend participation requirements such as zero-knowledge (ZK) identity verification or network buy-ins. Both of these implementations, as well as incentivizing an increasingly decentralized network, limit the possibility of Sybil attacks; Routing attacks are easily overcome by the encryption of all private data; and DDoS attacks can be limited by isolating excess network activity to each individual node.

To evaluate the validity of these statements we developed a prototype decentralized network  (Eden) that implements \texttt{Blossom} as its governing consensus protocol. The network is then tested across variable scales and against a variety of scenarios, as to provide a reasonable projection of performance. We ran our test across multiple AWS EC2 instances, with applications orchestrated within independent docker containers. We used Kubernetes to manage the orchestration of each docker container within EC2 instances \cite{burns2022kubernetes}. To benchmark our network we implemented the Prometheus Benchmarking tool to evaluate performance \cite{Prometheus_undated-lh}. 

With these results in mind, this paper outlines the contributions and improvements which have facilitated the aforementioned performance advancements in the following sections: a glossary of terms [\S\ref{term}], the discussion of related work [\S\ref{related}], our assumptions in developing \texttt{Blossom} [\S\ref{assum}], \texttt{Blossom}’s design [\S\ref{proto}], the projected performance of \texttt{Blossom}  [\S\ref{proj}]; system benchmarking [\S\ref{bench}], theoretical future additions to the protocol [\S\ref{future}], and a conclusion discussing our accomplishments, the implications of said accomplishments, and the prospective applications we believe will provide further improvements to our system [\S\ref{conclusion}].